\documentclass[12pt, a4paper]{article}
% --- 1. Cấu hình Tiếng Việt và Font chữ chuẩn ---
\usepackage[utf8]{inputenc}
\usepackage[T5]{fontenc}
\usepackage[vietnamese]{babel}
\usepackage{mathptmx} % Font Times New Roman đồng bộ cả chữ và công thức toán, không gây xung đột gói

\makeatletter
\renewcommand{\normalsize}{\@setfontsize\normalsize{13pt}{16pt}}
\makeatother
\normalsize

% --- 2. Gói Toán học & Ký hiệu bổ trợ ---
\usepackage{amsmath, amssymb, amsfonts, mathrsfs, amsthm}
\usepackage{graphicx}
\usepackage{fontawesome}
\usepackage{booktabs, tabularx, array, multirow}
\usepackage{enumitem}

% --- 3. Đồ họa TikZ, Bảng biến thiên & Hình không gian ---
\usepackage{tikz}
\usetikzlibrary{calc, intersections, angles, quotes, patterns, arrows.meta, 3d}
\usepackage{pgfplots}
\pgfplotsset{compat=1.18}
\usepackage{tkz-euclide}
\usepackage{tkz-tab}

\renewcommand{\vec}[1]{\overrightarrow{#1}}

% --- 4. Hộp sư phạm (Tcolorbox) ---
\usepackage[many]{tcolorbox}
\tcbuselibrary{skins, breakable}

% --- 5. Căn lề chuẩn văn bản giáo án ---
\usepackage{geometry}
\geometry{
    a4paper,
    left=25mm,
    right=15mm,
    top=20mm,
    bottom=20mm
}

% --- 6. Header / Footer chuẩn hóa ---
\usepackage{fancyhdr}
\usepackage{lastpage}
\pagestyle{fancy}
\fancyhf{}

\fancyhead[L]{\small \textit{Kế hoạch bài dạy Toán 11}}
\fancyhead[R]{\textbf{Trang \thepage/\pageref{LastPage}}}
\fancyfoot[L]{\small \textit{Tổ Toán - Tin}}
\fancyfoot[R]{\small \textit{Trường THCS\&THPT Hồng Vân}}
\renewcommand{\headrulewidth}{0.4pt}
\renewcommand{\footrulewidth}{0.4pt}

% --- 7. Giãn dòng & Giãn đoạn theo chuẩn Nghị định 30/2020/NĐ-CP ---
\usepackage{setspace}
\setstretch{1.15}
\setlength{\parindent}{1.0cm}
\setlength{\parskip}{5pt}

% Định nghĩa hộp sư phạm chuyên dụng
\newtcolorbox{pedabox}[2][]{
    enhanced,
    breakable,
    colback=blue!3!white,
    colframe=blue!65!black,
    fonttitle=\bfseries\small,
    title={#2},
    arc=2mm,
    boxrule=0.6pt,
    left=3mm, right=3mm, top=2mm, bottom=2mm,
    #1
}

\newtcolorbox{aibox}[2][]{
    enhanced,
    breakable,
    colback=teal!4!white,
    colframe=teal!70!black,
    fonttitle=\bfseries\small,
    title={#2},
    arc=2mm,
    boxrule=0.6pt,
    left=3mm, right=3mm, top=2mm, bottom=2mm,
    #1
}

\newtcolorbox{scaffoldbox}[2][]{
    enhanced,
    breakable,
    colback=orange!4!white,
    colframe=orange!80!black,
    fonttitle=\bfseries\small,
    title={#2},
    arc=2mm,
    boxrule=0.6pt,
    left=3mm, right=3mm, top=2mm, bottom=2mm,
    #1
}

\begin{document}

\begin{center}
    {\fontsize{14pt}{17pt}\selectfont \textbf{KẾ HOẠCH BÀI DẠY MÔN TOÁN LỚP 11}}\\[6pt]
    {\fontsize{16pt}{19pt}\selectfont \textbf{BÀI 2. CÔNG THỨC LƯỢNG GIÁC}}\\[4pt]
    {\fontsize{12pt}{15pt}\selectfont \textbf{Môn học: Toán 11} \ \textbar \ \textbf{Thời lượng: 2 tiết (Tiết 4 và Tiết 5 theo PPCT)}}
\end{center}

\vspace{5pt}

\section*{I. MỤC TIÊU}

\subsection*{1. Kiến thức, Năng lực}

\begin{itemize}[leftmargin=1.2cm]
    \item \textbf{Yêu cầu cần đạt (Nguyên văn CT GDPT 2018):}
    \begin{itemize}[leftmargin=0.6cm]
        \item Mô tả được các phép biến đổi lượng giác cơ bản: công thức cộng; công thức góc nhân đôi; công thức biến đổi tích thành tổng và công thức biến đổi tổng thành tích.
        \item Giải quyết được một số vấn đề thực tiễn gắn với giá trị lượng giác của góc lượng giác và các phép biến đổi lượng giác.
    \end{itemize}

    \item \textbf{Năng lực Toán học đặc thù:}
    \begin{itemize}[leftmargin=0.6cm]
        \item \textit{Năng lực tư duy và lập luận toán học:} So sánh và phân tích mối liên hệ logic giữa công thức cộng và công thức góc nhân đôi (khi $b = a$), giữa công thức cộng và công thức biến đổi tích thành tổng; lập luận suy diễn từ các hệ thức lượng giác để rút gọn biểu thức và chứng minh các đẳng thức lượng giác phức tạp.
        \item \textit{Năng lực mô hình hóa toán học:} Thiết lập mô hình toán học giải thích hiện tượng giao thoa sóng, sự cộng hưởng biên độ trong dao động điều hòa (Vật lí) khi tổng hợp hai dao động bằng công thức biến đổi tổng thành tích.
        \item \textit{Năng lực giải quyết vấn đề toán học:} Vận dụng linh hoạt các hệ thống công thức lượng giác để tính giá trị chính xác của các góc không đặc biệt (như $15^\circ, 75^\circ, \dfrac{\pi}{12}, \dfrac{5\pi}{12}$), rút gọn biểu thức và giải quyết các bài toán định lượng liên môn.
        \item \textit{Năng lực giao tiếp toán học:} Trình bày rõ ràng, mạch lạc các bước biến đổi lượng giác; ghi nhớ và đọc đúng các câu vè, mẹo sư phạm ghi nhớ công thức.
        \item \textit{Năng lực sử dụng công cụ, phương tiện học toán:} Khai thác máy tính cầm tay (MTCT) Casio fx-580VN X để kiểm tra nhanh kết quả rút gọn, phân giải biểu thức lượng giác và tính toán các giá trị số.
    \end{itemize}

    \item \textbf{Năng lực số (Theo Thông tư 02/2025/TT-BGDĐT):}
    \begin{itemize}[leftmargin=0.6cm]
        \item \textit{Mã 1.1NC1a:} Học sinh tra cứu, tìm kiếm và chọn lọc các bài thơ, mẹo ghi nhớ công thức cộng, công thức nhân đôi trên Internet hoặc các nền tảng học tập trực tuyến.
        \item \textit{Ứng dụng công cụ tính toán số:} Sử dụng MTCT Casio fx-580VN X với chức năng gán biến (\texttt{STO}) và tính giá trị biểu thức (\texttt{CALC}) để kiểm chứng tính đúng đắn của phép rút gọn biểu thức lượng giác.
    \end{itemize}

    \item \textbf{Năng lực AI (Theo Quyết định 2422/QĐ-BGDĐT):}
    \begin{itemize}[leftmargin=0.6cm]
        \item \textit{Mã 11.C2.1:} Học sinh trình bày được cách hệ thống AI tự động phân giải các biểu thức toán học, kiểm tra các bước chứng minh lượng giác và giải phương trình thông qua công cụ xử lý ngôn ngữ tự nhiên kết hợp toán học tính toán biểu trưng (Symbolic Math).
        \item \textit{Nhận diện giới hạn của AI:} Học sinh nhận thức được giới hạn của AI khi gặp các bài toán lượng giác có nhiều điều kiện ràng buộc hoặc dấu âm/dương ở từng góc phần tư cụ thể, từ đó biết cách phản biện và đối chiếu bài giải của AI.
    \end{itemize}

    \item \textbf{Năng lực chung:}
    \begin{itemize}[leftmargin=0.6cm]
        \item \textit{Tự chủ và tự học:} Tự đọc SGK, tự giác hoàn thành các nhiệm vụ học tập trên phiếu học tập cá nhân.
        \item \textit{Giao tiếp và hợp tác:} Tích cực thảo luận cặp đôi, nhóm nhỏ để chia sẻ phương pháp biến đổi công thức lượng giác tối ưu.
    \end{itemize}
\end{itemize}

\subsection*{2. Phẩm chất}

\begin{itemize}[leftmargin=1.2cm]
    \item \textbf{Chăm chỉ:} Kiên trì học thuộc hệ thống công thức lượng giác, tích cực làm bài tập rèn luyện kỹ năng biến đổi đại số lượng giác.
    \item \textbf{Trung thực:} Khách quan, trung thực trong quá trình tự làm bài, tự đối chiếu kết quả với máy tính cầm tay và AI.
    \item \textbf{Trách nhiệm:} Tích cực hỗ trợ bạn trong nhóm (đặc biệt học sinh trung bình -- yếu) cùng ghi nhớ công thức và tránh các sai lầm về dấu.
\end{itemize}

\section*{II. THIẾT BỊ DẠY HỌC VÀ HỌC LIỆU}

\begin{itemize}[leftmargin=1.2cm]
    \item \textbf{Giáo viên:}
    \begin{itemize}[leftmargin=0.6cm]
        \item Kế hoạch bài dạy chi tiết, bài trình chiếu PowerPoint trực quan hóa sơ đồ cây mối liên hệ giữa các công thức lượng giác.
        \item Máy chiếu, máy tính kết nối Internet, phần mềm giả lập MTCT Casio fx-580VN X.
        \item Bảng phụ / Giấy A0 in sẵn cây sơ đồ công thức lượng giác để học sinh thảo luận nhóm.
        \item Hệ thống Phiếu học tập phân tầng (Worksheet 1, Worksheet 2) và Rubric đánh giá năng lực 4 mức độ.
    \end{itemize}
    \item \textbf{Học sinh:}
    \begin{itemize}[leftmargin=0.6cm]
        \item SGK Toán 11, vở ghi bài, giấy nháp.
        \item Thước kẻ, bút màu, MTCT Casio fx-580VN X.
        \item Điện thoại thông minh (nếu có để thực hành tra cứu mẹo học tập và trải nghiệm prompt AI giải toán lượng giác).
    \end{itemize}
\end{itemize}

\section*{III. TIẾN TRÌNH DẠY HỌC}

% =========================================================================
% TIẾT 1 CỦA BÀI 2 (TIẾT 4 PPCT)
% =========================================================================
\subsection*{TIẾT 1. CÔNG THỨC CỘNG VÀ CÔNG THỨC GÓC NHÂN ĐÔI}
\textit{(Tiết 4 theo PPCT | Tuần 2)}

\subsubsection*{1. Hoạt động 1: Khởi động (Mở đầu)}
\textbf{a) Mục tiêu:} Tạo tình huống có vấn đề khi cần tính giá trị lượng giác của các góc không có sẵn trong bảng giá trị đặc biệt như $75^\circ, 15^\circ$, từ đó phát hiện sự sai lầm thường gặp: $\cos(a - b) \neq \cos a - \cos b$.

\textbf{b) Nội dung:} Giáo viên đưa ra câu hỏi:
\begin{itemize}[leftmargin=0.6cm]
    \item Ta đã biết giá trị lượng giác của các góc đặc biệt: $30^\circ, 45^\circ, 60^\circ, 90^\circ$. Làm thế nào để tính chính xác giá trị của $\cos 15^\circ$ mà không dùng máy tính?
    \item Một học sinh viết: $\cos 15^\circ = \cos(45^\circ - 30^\circ) = \cos 45^\circ - \cos 30^\circ = \dfrac{\sqrt{2}}{2} - \dfrac{\sqrt{3}}{2} = \dfrac{\sqrt{2} - \sqrt{3}}{2}$. Em hãy kiểm tra xem bạn tính đúng hay sai bằng máy tính cầm tay?
\end{itemize}

\textbf{c) Sản phẩm:} Học sinh kiểm tra bằng MTCT:
\begin{itemize}[leftmargin=0.6cm]
    \item $\dfrac{\sqrt{2} - \sqrt{3}}{2} \approx -0{,}1589 < 0$.
    \item Nhưng $15^\circ$ là góc nhọn (góc phần tư thứ I) nên $\cos 15^\circ > 0$, MTCT cho kết quả $\cos 15^\circ = \dfrac{\sqrt{6} + \sqrt{2}}{4} \approx 0{,}9659 > 0$.
    \item Kết luận: $\cos(a - b) \neq \cos a - \cos b$. Phép lấy cosin không có tính chất phân phối!
\end{itemize}

\textbf{d) Tổ chức thực hiện:}
\begin{itemize}[leftmargin=0.6cm]
    \item \textit{Chuyển giao nhiệm vụ:} Giáo viên chiếu bài toán lên bảng, yêu cầu cả lớp thực hành bấm MTCT kiểm tra trong 2 phút.
    \item \textit{Thực hiện nhiệm vụ (Giàn giáo sư phạm):}
    \begin{scaffoldbox}{Giàn giáo sư phạm 2.1: Cảnh báo sai lầm kinh điển}
        \textbf{Cảnh báo sai lầm cốt tử:}
        Tuyệt đối không áp dụng quy tắc nhân phân phối cho các hàm lượng giác!
        $$\cos(a \pm b) \neq \cos a \pm \cos b; \quad \sin(a \pm b) \neq \sin a \pm \sin b$$
        Cần phải có hệ thống công thức biến đổi riêng gọi là \textbf{Công thức cộng}.
    \end{scaffoldbox}
    \item \textit{Báo cáo, thảo luận:} 1 học sinh đứng tại chỗ trả lời kết quả so sánh trên MTCT.
    \item \textit{Kết luận, nhận định:} Giáo viên chốt lại mâu thuẫn nhận thức và dẫn dắt vào bài mới: Thiết lập công thức cộng và công thức nhân đôi.
\end{itemize}

\subsubsection*{2. Hoạt động 2: Hình thành kiến thức}

\paragraph{Nội dung 1: Công thức cộng}
\textbf{a) Mục tiêu:} Nhận biết, mô tả và ghi nhớ chính xác 6 công thức cộng đối với $\cos, \sin, \tan$; biết cách suy luận và áp dụng tính giá trị lượng giác.

\textbf{b) Nội dung:}
Với mọi góc lượng giác $a, b$, ta có:
\begin{align*}
    \cos(a - b) &= \cos a \cos b + \sin a \sin b \\
    \cos(a + b) &= \cos a \cos b - \sin a \sin b \\
    \sin(a - b) &= \sin a \cos b - \cos a \sin b \\
    \sin(a + b) &= \sin a \cos b + \cos a \sin b
\end{align*}
Với các góc $a, b$ thỏa mãn điều kiện để các biểu thức xác định:
\begin{align*}
    \tan(a - b) &= \dfrac{\tan a - \tan b}{1 + \tan a \tan b} \\
    \tan(a + b) &= \dfrac{\tan a + \tan b}{1 - \tan a \tan b}
\end{align*}

\textbf{c) Sản phẩm:} Học sinh ghi chép đầy đủ 6 công thức và vận dụng tính đúng $\cos 15^\circ$:
$$\cos 15^\circ = \cos(45^\circ - 30^\circ) = \cos 45^\circ \cos 30^\circ + \sin 45^\circ \sin 30^\circ = \dfrac{\sqrt{2}}{2} \cdot \dfrac{\sqrt{3}}{2} + \dfrac{\sqrt{2}}{2} \cdot \dfrac{1}{2} = \dfrac{\sqrt{6} + \sqrt{2}}{4}$$

\textbf{d) Tổ chức thực hiện (Tích hợp Năng lực số 1.1NC1a & Giàn giáo sư phạm):}
\begin{itemize}[leftmargin=0.6cm]
    \item \textit{Chuyển giao nhiệm vụ:} Giáo viên trình bày công thức, yêu cầu học sinh thảo luận cặp đôi tìm cách ghi nhớ nhanh các công thức cộng.
    \item \textit{Thực hiện nhiệm vụ (Mẹo sư phạm cho HS TB-Yếu):}
    \begin{scaffoldbox}{Giàn giáo sư phạm 2.2: Câu vè dân gian ghi nhớ công thức cộng}
        \textbf{Bài vè truyền miệng nổi tiếng:}\\
        \textit{''Cos thì cos cos sin sin,}\\
        \textit{Sin thì sin cos cos sin rõ ràng.}\\
        \textit{Cos thì đổi dấu chớ màng,}\\
        \textit{Sin thì giữ dấu đàng hoàng không sai.''}\\
        \textbf{Giải thích quy tắc dấu:}\\
        - Với $\cos$: dấu trong ngoặc là trừ ($-$) thì vế phải là cộng ($+$); dấu trong ngoặc là cộng ($+$) thì vế phải là trừ ($-$).\\
        - Với $\sin$: dấu trong ngoặc là trừ ($-$) thì vế phải giữ nguyên trừ ($-$), cộng ($+$) thì vế phải là cộng ($+$).
    \end{scaffoldbox}
    \item \textit{Báo cáo, thảo luận:} Cả lớp đọc đồng thanh câu vè ghi nhớ. 1 học sinh đại diện áp dụng tính nhanh $\sin 75^\circ = \sin(45^\circ + 30^\circ) = \dfrac{\sqrt{6} + \sqrt{2}}{4}$.
    \item \textit{Kết luận, nhận định:} Giáo viên chuẩn hóa công thức, lưu ý điều kiện xác định của công thức $\tan$.
\end{itemize}

\paragraph{Nội dung 2: Công thức góc nhân đôi và công thức hạ bậc}
\textbf{a) Mục tiêu:} Suy luận được công thức góc nhân đôi từ công thức cộng khi đặt $b = a$; nắm vững 3 dạng biểu diễn của $\cos 2a$; viết được công thức hạ bậc.

\textbf{b) Nội dung:}
\begin{itemize}[leftmargin=0.6cm]
    \item Trong các công thức cộng $\sin(a + b), \cos(a + b), \tan(a + b)$, thay $b = a$, ta thu được \textbf{Công thức góc nhân đôi}:
    \begin{align*}
        \sin 2a &= 2\sin a \cos a \\
        \cos 2a &= \cos^2 a - \sin^2 a = 2\cos^2 a - 1 = 1 - 2\sin^2 a \\
        \tan 2a &= \dfrac{2\tan a}{1 - \tan^2 a} \quad (\text{với các điều kiện xác định})
    \end{align*}
    \item Từ công thức $\cos 2a = 2\cos^2 a - 1 = 1 - 2\sin^2 a$, ta suy ra \textbf{Công thức hạ bậc}:
    $$\cos^2 a = \dfrac{1 + \cos 2a}{2}; \quad \sin^2 a = \dfrac{1 - \cos 2a}{2}$$
\end{itemize}

\textbf{c) Sản phẩm:} Học sinh tự suy luận ra công thức góc nhân đôi vào vở và vận dụng tính $\sin 2a, \cos 2a$ khi biết $\sin a = \dfrac{3}{5}$ ($0 < a < \dfrac{\pi}{2}$).

\textbf{d) Tổ chức thực hiện:}
\begin{itemize}[leftmargin=0.6cm]
    \item \textit{Chuyển giao nhiệm vụ:} Giáo viên giao bài tập: Cho $b = a$ trong công thức cộng, hãy viết ra các công thức thu được.
    \item \textit{Thực hiện nhiệm vụ:} Học sinh làm việc độc lập 3 phút, sau đó đổi vở kiểm tra.
    \item \textit{Báo cáo, thảo luận:} 1 học sinh lên bảng trình bày 3 biến thể của $\cos 2a$.
    \item \textit{Kết luận, nhận định:} Giáo viên nhấn mạnh: Góc ở vế trái là $2a$ thì góc ở vế phải là $a$ (giảm đi một nửa). Mở rộng: $\sin 4a = 2\sin 2a \cos 2a; \cos a = 2\cos^2\dfrac{a}{2} - 1$.
\end{itemize}

\subsubsection*{3. Hoạt động 3: Luyện tập}
\textbf{a) Mục tiêu:} Rèn luyện kỹ năng tính giá trị lượng giác, rút gọn biểu thức bằng công thức cộng và công thức nhân đôi; kiểm tra kết quả bằng MTCT.

\textbf{b) Nội dung:} Học sinh hoàn thành Phiếu học tập số 1 (Worksheet 1):
\begin{itemize}[leftmargin=0.6cm]
    \item \textbf{Bài 1:} Không dùng máy tính cầm tay, hãy tính giá trị của:
    $$A = \sin 47^\circ \cos 17^\circ - \cos 47^\circ \sin 17^\circ; \quad B = \cos\dfrac{7\pi}{12} \cos\dfrac{\pi}{12} + \sin\dfrac{7\pi}{12} \sin\dfrac{\pi}{12}$$
    \item \textbf{Bài 2:} Cho $\cos a = -\dfrac{4}{5}$ với $\pi < a < \dfrac{3\pi}{2}$. Tính $\sin 2a, \cos 2a, \tan 2a$.
    \item \textbf{Bài 3:} Rút gọn biểu thức: $P = \dfrac{\sin 2x}{2\cos x} - \sin x$.
\end{itemize}

\textbf{c) Sản phẩm:} Lời giải chi tiết:
\begin{itemize}[leftmargin=0.6cm]
    \item \textbf{Bài 1:}
    \begin{itemize}[leftmargin=0.5cm]
        \item $A = \sin(47^\circ - 17^\circ) = \sin 30^\circ = \dfrac{1}{2}$.
        \item $B = \cos\left(\dfrac{7\pi}{12} - \dfrac{\pi}{12}\right) = \cos\dfrac{6\pi}{12} = \cos\dfrac{\pi}{2} = 0$.
    \end{itemize}
    \item \textbf{Bài 2:}
    \begin{itemize}[leftmargin=0.5cm]
        \item Vì $\pi < a < \dfrac{3\pi}{2}$ (góc phần tư thứ III) nên $\sin a < 0$.\\
        $\sin a = -\sqrt{1 - \cos^2 a} = -\sqrt{1 - \left(-\dfrac{4}{5}\right)^2} = -\dfrac{3}{5}$.
        \item $\sin 2a = 2\sin a \cos a = 2 \cdot \left(-\dfrac{3}{5}\right) \cdot \left(-\dfrac{4}{5}\right) = \dfrac{24}{25}$.
        \item $\cos 2a = 2\cos^2 a - 1 = 2 \cdot \left(-\dfrac{4}{5}\right)^2 - 1 = \dfrac{32}{25} - 1 = \dfrac{7}{25}$.
        \item $\tan 2a = \dfrac{\sin 2a}{\cos 2a} = \dfrac{\dfrac{24}{25}}{\dfrac{7}{25}} = \dfrac{24}{7}$.
    \end{itemize}
    \item \textbf{Bài 3:} Ta có $\sin 2x = 2\sin x \cos x$, do đó:
    $$P = \dfrac{2\sin x \cos x}{2\cos x} - \sin x = \sin x - \sin x = 0$$
\end{itemize}

\textbf{d) Tổ chức thực hiện:}
\begin{itemize}[leftmargin=0.6cm]
    \item \textit{Chuyển giao nhiệm vụ:} Học sinh làm bài theo cặp đôi trong 7 phút.
    \item \textit{Thực hiện nhiệm vụ (Hỗ trợ HS TB-Yếu):} Giáo viên nhắc nhở học sinh yếu xác định đúng dấu của $\sin a$ khi khai căn ở Bài 2.
    \item \textit{Báo cáo, thảo luận:} 3 học sinh lên bảng trình bày 3 bài toán. Cả lớp đối chiếu kết quả bấm MTCT.
    \item \textit{Kết luận, nhận định:} Giáo viên nhận xét, nhấn mạnh kỹ thuật nhận dạng công thức theo chiều ngược lại (từ vế phải sang vế trái).
\end{itemize}

\subsubsection*{4. Hoạt động 4: Vận dụng}
\textbf{a) Mục tiêu:} Vận dụng công thức cộng và nhân đôi vào bài toán thực tiễn tính toán tầm xa của vật ném xiên trong Vật lí.

\textbf{b) Nội dung:}
Một quả bóng được đá lên từ mặt đất với vận tốc ban đầu $v_0$ và góc phóng $\alpha$ so với phương ngang. Tầm xa $L$ của quả bóng được tính theo công thức Vật lí:
$$L = \dfrac{2 v_0^2 \sin\alpha \cos\alpha}{g}$$
trong đó $g$ là gia tốc trọng trường ($g \approx 9{,}8\text{ m/s}^2$).
\begin{enumerate}[leftmargin=0.5cm]
    \item Hãy áp dụng công thức góc nhân đôi để viết lại công thức tính tầm xa $L$ gọn hơn.
    \item Với vận tốc ban đầu $v_0$ không đổi, góc phóng $\alpha$ bằng bao nhiêu độ thì quả bóng đạt tầm xa cực đại?
\end{enumerate}

\textbf{c) Sản phẩm:} Lời giải:
\begin{enumerate}[leftmargin=0.5cm]
    \item Áp dụng công thức $2\sin\alpha \cos\alpha = \sin 2\alpha$, ta có:
    $$L = \dfrac{v_0^2 \sin 2\alpha}{g}$$
    \item Tầm xa $L$ lớn nhất khi $\sin 2\alpha$ đạt giá trị lớn nhất.\\
    Vì $\sin 2\alpha \le 1$ nên $\max L = \dfrac{v_0^2}{g} \iff \sin 2\alpha = 1 \iff 2\alpha = 90^\circ \iff \alpha = 45^\circ$.\\
    Vậy để đạt tầm xa cực đại, góc sút phải là $45^\circ$.
\end{enumerate}

\textbf{d) Tổ chức thực hiện:}
\begin{itemize}[leftmargin=0.6cm]
    \item \textit{Chuyển giao nhiệm vụ:} Giáo viên giao bài toán liên môn Vật lý -- Thể thao cho các nhóm 4 người.
    \item \textit{Thực hiện nhiệm vụ:} Các nhóm thảo luận giải quyết nhanh trong 3 phút.
    \item \textit{Báo cáo, thảo luận:} 1 nhóm xung phong trình bày, giải thích góc $45^\circ$ trong thực tế thi đấu bóng đá/ném lao.
    \item \textit{Kết luận, nhận định:} Giáo viên tổng kết Tiết 1 của Bài 2, dặn dò học sinh ghi nhớ các công thức cộng và nhân đôi.
\end{itemize}

% =========================================================================
% TIẾT 2 CỦA BÀI 2 (TIẾT 5 PPCT)
% =========================================================================
\vspace{10pt}
\subsection*{TIẾT 2. CÔNG THỨC BIẾN ĐỔI TÍCH THÀNH TỔNG VÀ TỔNG THÀNH TÍCH}
\textit{(Tiết 5 theo PPCT | Tuần 2)}

\subsubsection*{1. Hoạt động 1: Khởi động (Mở đầu)}
\textbf{a) Mục tiêu:} Tạo tình huống cần tính toán tổng hai dao động điều hòa hoặc tích các giá trị lượng giác, dẫn dắt vào nhu cầu biến đổi giữa tích và tổng.

\textbf{b) Nội dung:}
Cho hai âm thanh phát ra cùng lúc với tần số gần nhau, tạo ra dao động sóng âm có biểu thức:
$$y = \sin(100\pi t) + \sin(104\pi t)$$
Làm thế nào để đưa biểu thức trên về dạng tích của một hàm biên độ biến thiên chậm và một sóng hình sin để phân tích hiện tượng ''phách'' (sự biến thiên tuần hoàn của âm lượng)?

\textbf{c) Sản phẩm:} Học sinh nhận thức được: Cần có công thức biến đổi tổng $\sin u + \sin v$ thành tích của các hàm lượng giác.

\textbf{d) Tổ chức thực hiện:}
\begin{itemize}[leftmargin=0.6cm]
    \item \textit{Chuyển giao nhiệm vụ:} Giáo viên giới thiệu ngắn gọn hiện tượng âm thanh thực tế, nêu biểu thức.
    \item \textit{Thực hiện nhiệm vụ:} Học sinh quan sát, nhận thấy chưa thể cộng trực tiếp hai góc khác nhau.
    \item \textit{Báo cáo, thảo luận:} Học sinh nêu nhu cầu: Cần công thức chuyển tổng thành tích.
    \item \textit{Kết luận, nhận định:} Giáo viên dẫn dắt vào bài học mới.
\end{itemize}

\subsubsection*{2. Hoạt động 2: Hình thành kiến thức}

\paragraph{Nội dung 1: Công thức biến đổi tích thành tổng}
\textbf{a) Mục tiêu:} Suy luận được công thức biến đổi tích thành tổng bằng cách cộng/trừ từng vế các công thức cộng; ghi nhớ 3 công thức cơ bản.

\textbf{b) Nội dung:}
\begin{itemize}[leftmargin=0.6cm]
    \item Xuất phát từ:
    \begin{align*}
        \cos(a - b) &= \cos a \cos b + \sin a \sin b \\
        \cos(a + b) &= \cos a \cos b - \sin a \sin b
    \end{align*}
    Cộng từng vế: $\cos(a - b) + \cos(a + b) = 2\cos a \cos b \implies \cos a \cos b = \dfrac{1}{2}[\cos(a - b) + \cos(a + b)]$.
    \item Trừ từng vế: $\cos(a - b) - \cos(a + b) = 2\sin a \sin b \implies \sin a \sin b = \dfrac{1}{2}[\cos(a - b) - \cos(a + b)]$.
    \item Tương tự từ công thức $\sin(a \pm b)$:
    $$\sin a \cos b = \dfrac{1}{2}[\sin(a - b) + \sin(a + b)]$$
\end{itemize}

\textbf{c) Sản phẩm:} Học sinh ghi nhớ 3 công thức biến đổi tích thành tổng:
\begin{align*}
    \cos a \cos b &= \dfrac{1}{2}[\cos(a - b) + \cos(a + b)] \\
    \sin a \sin b &= \dfrac{1}{2}[\cos(a - b) - \cos(a + b)] \\
    \sin a \cos b &= \dfrac{1}{2}[\sin(a - b) + \sin(a + b)]
\end{align*}

\textbf{d) Tổ chức thực hiện:}
\begin{itemize}[leftmargin=0.6cm]
    \item \textit{Chuyển giao nhiệm vụ:} Giáo viên hướng dẫn học sinh cộng và trừ từng vế của các công thức cộng $\cos(a \pm b)$ để rút ra $\cos a \cos b$ và $\sin a \sin b$.
    \item \textit{Thực hiện nhiệm vụ (Giàn giáo cho HS TB-Yếu):}
    \begin{scaffoldbox}{Giàn giáo sư phạm 2.3: Chú ý dấu ở tích sin nhân sin}
        \textbf{Lưu ý đặc biệt:} Công thức $\sin a \sin b$ rất hay bị học sinh viết nhầm dấu!
        Nhớ rằng vế phải là: $\dfrac{1}{2}[\cos(a - b) \mathbf{-} \cos(a + b)]$. Luôn luôn lấy góc hiệu $(a - b)$ đứng trước!
    \end{scaffoldbox}
    \item \textit{Báo cáo, thảo luận:} 1 học sinh lên bảng trình bày suy luận cho công thức $\sin a \cos b$.
    \item \textit{Kết luận, nhận định:} Giáo viên chuẩn hóa công thức, yêu cầu học sinh đóng khung.
\end{itemize}

\paragraph{Nội dung 2: Công thức biến đổi tổng thành tích}
\textbf{a) Mục tiêu:} Suy luận công thức biến đổi tổng thành tích bằng cách đặt ẩn phụ $u = a + b, v = a - b$; ghi nhớ 4 công thức cơ bản và bài thơ ghi nhớ.

\textbf{b) Nội dung:}
\begin{itemize}[leftmargin=0.6cm]
    \item Đặt $u = a + b$ và $v = a - b \implies a = \dfrac{u + v}{2}$ và $b = \dfrac{u - v}{2}$.
    \item Thay vào các công thức tích thành tổng, ta thu được \textbf{Công thức biến đổi tổng thành tích}:
    \begin{align*}
        \cos u + \cos v &= 2\cos\dfrac{u + v}{2} \cos\dfrac{u - v}{2} \\
        \cos u - \cos v &= -2\sin\dfrac{u + v}{2} \sin\dfrac{u - v}{2} \\
        \sin u + \sin v &= 2\sin\dfrac{u + v}{2} \cos\dfrac{u - v}{2} \\
        \sin u - \sin v &= 2\cos\dfrac{u + v}{2} \sin\dfrac{u - v}{2}
    \end{align*}
\end{itemize}

\textbf{c) Sản phẩm:} Học sinh thuộc 4 công thức và áp dụng giải bài toán khởi động:
$$\sin(100\pi t) + \sin(104\pi t) = 2\sin\left(\dfrac{100\pi t + 104\pi t}{2}\right) \cos\left(\dfrac{100\pi t - 104\pi t}{2}\right) = 2\cos(2\pi t) \sin(102\pi t)$$

\textbf{d) Tổ chức thực hiện (Tích hợp Giàn giáo sư phạm & Năng lực AI):}
\begin{itemize}[leftmargin=0.6cm]
    \item \textit{Chuyển giao nhiệm vụ:} Giáo viên yêu cầu học sinh ghi nhớ 4 công thức và giải bài toán âm thanh.
    \item \textit{Thực hiện nhiệm vụ (Mẹo sư phạm ghi nhớ):}
    \begin{scaffoldbox}{Giàn giáo sư phạm 2.4: Thơ ghi nhớ tổng thành tích}
        \textbf{Bài vè truyền miệng:}\\
        \textit{''Cos cộng cos bằng hai cos cos,}\\
        \textit{Cos trừ cos bằng trừ hai sin sin.}\\
        \textit{Sin cộng sin bằng hai sin cos,}\\
        \textit{Sin trừ sin bằng hai cos sin.''}
    \end{scaffoldbox}
    \item \textit{Tích hợp Năng lực AI (Mã 11.C2.1):}
    \begin{aibox}{Thực hành Năng lực AI: Kiểm tra phân giải biểu thức lượng giác}
        \textbf{Kỹ thuật Prompt hỏi AI:}\\
        \texttt{"Hãy phân giải biểu thức $\sin(5x) + \sin(3x)$ thành tích và giải thích chi tiết các bước áp dụng công thức biến đổi tổng thành tích."}\\
        $\implies$ Học sinh đối chiếu lời giải của AI: $u = 5x, v = 3x \implies \dfrac{u + v}{2} = 4x, \dfrac{u - v}{2} = x \implies 2\sin(4x)\cos(x)$. Nhận xét: AI tính toán rất nhanh và chính xác các biểu thức đại số giải tích chuẩn mực.
    \end{aibox}
    \item \textit{Báo cáo, thảo luận:} Học sinh đọc đồng thanh bài vè, 1 học sinh giải thích kết quả AI.
    \item \textit{Kết luận, nhận định:} Giáo viên nhấn mạnh: Cần kiểm tra kỹ các hệ số $\dfrac{u + v}{2}$ và $\dfrac{u - v}{2}$, tránh quên chia 2.
\end{itemize}

\subsubsection*{3. Hoạt động 3: Luyện tập}
\textbf{a) Mục tiêu:} Rèn luyện thành thạo kỹ năng biến đổi qua lại giữa tích và tổng; rút gọn biểu thức phân thức lượng giác.

\textbf{b) Nội dung:} Học sinh hoàn thành Phiếu học tập số 2 (Worksheet 2):
\begin{itemize}[leftmargin=0.6cm]
    \item \textbf{Bài 1:} Biến đổi thành tổng:
    $$A = 4\cos(3x) \cos(2x); \quad B = \sin\left(x + \dfrac{\pi}{4}\right) \sin\left(x - \dfrac{\pi}{4}\right)$$
    \item \textbf{Bài 2:} Biến đổi thành tích:
    $$C = \cos(4x) + \cos(2x); \quad D = 1 + 2\cos x$$
    \item \textbf{Bài 3:} Rút gọn biểu thức lượng giác:
    $$M = \dfrac{\sin x + \sin 3x + \sin 5x}{\cos x + \cos 3x + \cos 5x}$$
\end{itemize}

\textbf{c) Sản phẩm:} Lời giải chi tiết:
\begin{itemize}[leftmargin=0.6cm]
    \item \textbf{Bài 1:}
    \begin{itemize}[leftmargin=0.5cm]
        \item $A = 4 \cdot \dfrac{1}{2}[\cos(3x - 2x) + \cos(3x + 2x)] = 2(\cos x + \cos 5x) = 2\cos x + 2\cos 5x$.
        \item $B = \dfrac{1}{2}\left[\cos\left(\left(x + \dfrac{\pi}{4}\right) - \left(x - \dfrac{\pi}{4}\right)\right) - \cos\left(\left(x + \dfrac{\pi}{4}\right) + \left(x - \dfrac{\pi}{4}\right)\right)\right] = \dfrac{1}{2}\left[\cos\dfrac{\pi}{2} - \cos 2x\right] = -\dfrac{1}{2}\cos 2x$.
    \end{itemize}
    \item \textbf{Bài 2:}
    \begin{itemize}[leftmargin=0.5cm]
        \item $C = 2\cos\left(\dfrac{4x + 2x}{2}\right) \cos\left(\dfrac{4x - 2x}{2}\right) = 2\cos(3x)\cos x$.
        \item $D = 2\left(\dfrac{1}{2} + \cos x\right) = 2\left(\cos\dfrac{\pi}{3} + \cos x\right) = 2 \cdot 2\cos\left(\dfrac{x + \dfrac{\pi}{3}}{2}\right) \cos\left(\dfrac{x - \dfrac{\pi}{3}}{2}\right) = 4\cos\left(\dfrac{x}{2} + \dfrac{\pi}{6}\right) \cos\left(\dfrac{x}{2} - \dfrac{\pi}{6}\right)$.
    \end{itemize}
    \item \textbf{Bài 3:}
    Nhóm số hạng thứ nhất và thứ ba:
    \begin{align*}
        \text{Tử số} &= (\sin 5x + \sin x) + \sin 3x = 2\sin 3x \cos 2x + \sin 3x = \sin 3x (2\cos 2x + 1) \\
        \text{Mẫu số} &= (\cos 5x + \cos x) + \cos 3x = 2\cos 3x \cos 2x + \cos 3x = \cos 3x (2\cos 2x + 1)
    \end{align*}
    Do đó:
    $$M = \dfrac{\sin 3x (2\cos 2x + 1)}{\cos 3x (2\cos 2x + 1)} = \dfrac{\sin 3x}{\cos 3x} = \tan 3x$$
\end{itemize}

\textbf{d) Tổ chức thực hiện:}
\begin{itemize}[leftmargin=0.6cm]
    \item \textit{Chuyển giao nhiệm vụ:} Học sinh làm bài độc lập trong 8 phút.
    \item \textit{Thực hiện nhiệm vụ (Giàn giáo cho HS TB-Yếu):}
    \begin{scaffoldbox}{Giàn giáo sư phạm 2.5: Chiến thuật nhóm ghép đối xứng}
        Khi gặp tổng 3 hàm lượng giác như $\sin x + \sin 3x + \sin 5x$, ta nên ghép số hạng góc nhỏ nhất ($x$) với góc lớn nhất ($5x$) vì $\dfrac{x + 5x}{2} = 3x$, từ đó làm xuất hiện nhân tử chung $\sin 3x$.
    \end{scaffoldbox}
    \item \textit{Báo cáo, thảo luận:} 3 học sinh lên bảng trình bày, cả lớp nhận xét.
    \item \textit{Kết luận, nhận định:} Giáo viên chuẩn hóa kỹ thuật ghép đối xứng và kỹ thuật đưa số thực về giá trị lượng giác để áp dụng công thức tổng thành tích.
\end{itemize}

\subsubsection*{4. Hoạt động 4: Vận dụng}
\textbf{a) Mục tiêu:} Vận dụng công thức tổng thành tích để giải quyết bài toán giao thoa hai nguồn sóng kết hợp trong Vật lí.

\textbf{b) Nội dung:}
Hai nguồn phát sóng kết hợp $S_1, S_2$ dao động cùng pha, cùng biên độ $A = 2\text{ cm}$ và tần số $f = 50\text{ Hz}$. Tại một điểm $M$ trong môi trường, hai sóng truyền tới tạo nên hai dao động:
$$u_1 = 2\cos\left(100\pi t - \dfrac{2\pi d_1}{\lambda}\right) \text{ (cm)}; \quad u_2 = 2\cos\left(100\pi t - \dfrac{2\pi d_2}{\lambda}\right) \text{ (cm)}$$
\begin{enumerate}[leftmargin=0.5cm]
    \item Sử dụng công thức biến đổi tổng thành tích, hãy xác định phương trình dao động tổng hợp $u = u_1 + u_2$ tại điểm $M$.
    \item Biên độ dao động tổng hợp tại $M$ phụ thuộc vào hiệu đường đi $(d_2 - d_1)$ như thế nào? Điểm $M$ dao động với biên độ cực đại khi nào?
\end{enumerate}

\textbf{c) Sản phẩm:} Lời giải:
\begin{enumerate}[leftmargin=0.5cm]
    \item Áp dụng công thức $\cos u + \cos v = 2\cos\dfrac{u + v}{2} \cos\dfrac{u - v}{2}$:
    $$u = u_1 + u_2 = 2 \cdot 2 \cos\left(\dfrac{\dfrac{2\pi(d_2 - d_1)}{\lambda}}{2}\right) \cos\left(100\pi t - \dfrac{\pi(d_1 + d_2)}{\lambda}\right)$$
    $$u = 4\cos\left(\dfrac{\pi(d_2 - d_1)}{\lambda}\right) \cos\left(100\pi t - \dfrac{\pi(d_1 + d_2)}{\lambda}\right) \text{ (cm)}$$
    \item Biên độ dao động tại $M$ là: $A_M = 4 \left|\cos\left(\dfrac{\pi(d_2 - d_1)}{\lambda}\right)\right|$.\\
    Biên độ cực đại $A_{\max} = 4\text{ cm} \iff \left|\cos\left(\dfrac{\pi(d_2 - d_1)}{\lambda}\right)\right| = 1 \iff \dfrac{\pi(d_2 - d_1)}{\lambda} = k\pi \iff d_2 - d_1 = k\lambda$ ($k \in \mathbb{Z}$).\\
    (Hiệu đường đi bằng một số nguyên lần bước sóng).
\end{enumerate}

\textbf{d) Tổ chức thực hiện:}
\begin{itemize}[leftmargin=0.6cm]
    \item \textit{Chuyển giao nhiệm vụ:} Giáo viên giao bài toán liên môn Vật lý 11 (Hiện tượng giao thoa sóng) cho các nhóm học tập.
    \item \textit{Thực hiện nhiệm vụ:} Nhóm thảo luận trong 4 phút, thực hiện phép cộng lượng giác.
    \item \textit{Báo cáo, thảo luận:} Đại diện nhóm chuyên Lí phát biểu giải thích mối liên kết tuyệt vời giữa Toán học và hiện tượng sóng nước giao thoa.
    \item \textit{Kết luận, nhận định:} Giáo viên tổng kết toàn bộ Bài 2: Hệ thống 4 nhóm công thức (Cộng, Nhân đôi, Tích thành tổng, Tổng thành tích) là bộ công cụ quyền năng nhất của lượng giác. Hướng dẫn học sinh hoàn thành sơ đồ tư duy tóm tắt bài học.
\end{itemize}

% =========================================================================
% PHỤ LỤC
% =========================================================================
\newpage
\section*{IV. PHỤ LỤC HỌC LIỆU BỔ TRỢ}

\subsection*{1. Hệ thống Phiếu học tập (Worksheet) phân tầng bậc thang}

\begin{pedabox}{PHIẾU HỌC TẬP SỐ 1: CÔNG THỨC CỘNG VÀ NHÂN ĐÔI (Phân tầng 3 mức độ)}
\textbf{Họ và tên học sinh:} \dotfill \textbf{Lớp:} 11A\dots \ \textbf{Ngày:} \dotfill

\textbf{MỨC 1: NHẬN BIẾT (Điền khuyết có giàn giáo)}
\begin{enumerate}[leftmargin=0.5cm]
    \item Điền dấu ($+$ hoặc $-$) thích hợp vào chỗ trống:\\
    a) $\cos(a + b) = \cos a \cos b \dots\dots \sin a \sin b$; \quad b) $\sin(a - b) = \sin a \cos b \dots\dots \cos a \sin b$.
    \item Điền công thức nhân đôi:\\
    $\sin 2a = \dots\dots\dots\dots$; \quad $\cos 2a = \cos^2 a - \dots\dots\dots\dots = 2\cos^2 a - \dots\dots$
\end{enumerate}

\textbf{MỨC 2: THÔNG HIỂU (Áp dụng công thức tính giá trị)}
\begin{enumerate}[leftmargin=0.5cm]
    \item Tính giá trị chính xác của $\tan 105^\circ$ bằng cách viết $105^\circ = 60^\circ + 45^\circ$.\\
    \textit{Giàn giáo gợi ý:} Sử dụng công thức $\tan(a + b) = \dfrac{\tan a + \tan b}{1 - \tan a \tan b}$ với $a = 60^\circ, b = 45^\circ$.
    \item Cho $\sin x = \dfrac{1}{3}$ với $0 < x < \dfrac{\pi}{2}$. Tính $\cos 2x$ và $\sin 2x$.
\end{enumerate}

\textbf{MỨC 3: VẬN DỤNG (Chứng minh đẳng thức)}
\begin{enumerate}[leftmargin=0.5cm]
    \item Chứng minh rằng với mọi góc $x$ làm cho biểu thức có nghĩa, ta có:
    $$\dfrac{1 - \cos 2x + \sin 2x}{1 + \cos 2x + \sin 2x} = \tan x$$
\end{enumerate}
\end{pedabox}

\vspace{8pt}

\begin{pedabox}{PHIẾU HỌC TẬP SỐ 2: BIẾN ĐỔI TÍCH THÀNH TỔNG VÀ TỔNG THÀNH TÍCH}
\textbf{MỨC 1: NHẬN BIẾT}
\begin{enumerate}[leftmargin=0.5cm]
    \item Điền vào chỗ trống theo bài vè:\\
    a) $\cos a \cos b = \dfrac{1}{2}[\dots\dots\dots\dots + \cos(a + b)]$; \quad b) $\sin u + \sin v = 2\sin\dfrac{u+v}{2}\dots\dots\dots\dots$
    \item Đúng (Đ) hay Sai (S):\\
    a) $\cos u - \cos v = 2\sin\dfrac{u+v}{2}\sin\dfrac{u-v}{2}$ (\dots) \quad b) $\sin a \cos b = \dfrac{1}{2}[\sin(a-b) + \sin(a+b)]$ (\dots)
\end{enumerate}

\textbf{MỨC 2: THÔNG HIỂU}
\begin{enumerate}[leftmargin=0.5cm]
    \item Tính giá trị biểu thức không dùng máy tính: $P = \cos 75^\circ \sin 15^\circ$.\\
    \textit{Giàn giáo:} Dùng công thức $\sin b \cos a = \dfrac{1}{2}[\sin(b - a) + \sin(b + a)]$ với $b = 15^\circ, a = 75^\circ$.
    \item Rút gọn biểu thức: $Q = \cos\left(x + \dfrac{\pi}{3}\right) + \cos\left(x - \dfrac{\pi}{3}\right)$.
\end{enumerate}

\textbf{MỨC 3: VẬN DỤNG}
\begin{enumerate}[leftmargin=0.5cm]
    \item Rút gọn biểu thức phân thức lượng giác chứa 4 số hạng:
    $$K = \dfrac{\sin x + \sin 2x + \sin 3x + \sin 4x}{\cos x + \cos 2x + \cos 3x + \cos 4x}$$
    \textit{Giàn giáo:} Nhóm $(\sin 4x + \sin x)$ và $(\sin 3x + \sin 2x)$ để thấy xuất hiện góc chung $\dfrac{5x}{2}$.
\end{enumerate}
\end{pedabox}

\subsection*{2. Rubric đánh giá năng lực học sinh theo 4 mức độ}

\begin{table}[h!]
\centering
\small
\begin{tabularx}{\textwidth}{|p{2.8cm}|X|X|X|X|}
\hline
\textbf{Tiêu chí} & \textbf{Chưa đạt (Mức 1)} & \textbf{Đạt (Mức 2)} & \textbf{Khá (Mức 3)} & \textbf{Tốt (Mức 4)} \\
\hline
\textbf{1. Công thức cộng và nhân đôi} & Nhầm lẫn phép lấy hàm lượng giác với phép nhân phân phối; sai dấu $\cos(a \pm b)$. & Thuộc bài thơ ghi nhớ; tính được giá trị góc đặc biệt dạng $45^\circ \pm 30^\circ$. & Áp dụng thành thạo công thức nhân đôi, hạ bậc; tính đúng khi góc ở góc phần tư II, III. & Biến đổi linh hoạt hai chiều công thức; chứng minh được các hệ thức lượng giác phức tạp. \\
\hline
\textbf{2. Biến đổi tích -- tổng} & Nhầm dấu công thức $\sin a \sin b$; quên chia đôi góc trong tổng thành tích. & Nhớ công thức cơ bản; biến đổi được tích hai góc cụ thể thành tổng. & Biết chiến thuật ghép đối xứng góc lớn với góc nhỏ; rút gọn được phân thức bậc 3. & Phân tích và xử lý trọn vẹn các biểu thức nhiều số hạng đối xứng; ứng dụng giải phương trình. \\
\hline
\textbf{3. Vận dụng liên môn (STEM)} & Chưa nhận ra vai trò của công thức lượng giác trong Vật lí. & Thay số vào công thức cho trước để tính toán đại lượng. & Hiểu được bản chất cực đại của tầm ném xiên và biên độ sóng giao thoa qua hàm $\sin, \cos$. & Tự thiết lập mô hình lượng giác để mô tả và tối ưu hóa hiện tượng vật lý trong thực tế. \\
\hline
\textbf{4. Năng lực số \& Năng lực AI} & Chưa biết tra cứu mẹo học tập; không biết dùng MTCT kiểm tra biểu thức. & Biết dùng phím \texttt{CALC} trên MTCT kiểm tra giá trị rút gọn của biểu thức. & Đặt prompt chuẩn cho AI phân giải biểu thức; kiểm chứng được các bước giải. & Phân tích được thuật ngữ toán học biểu trưng của AI; phát hiện lỗi logic hoặc điều kiện thiếu của AI. \\
\hline
\end{tabularx}
\end{table}

\end{document}